\documentclass[12pt]{article}
\usepackage{amsmath,amssymb,amsthm}
\usepackage{array}
\usepackage{boxedminipage}
\usepackage[margin=1in]{geometry}
\usepackage{fancyhdr}
\usepackage{ragged2e}

\newenvironment{problem}[1][]{%
  \bigskip% put space before problem statement box %
  \noindent\begin{boxedminipage}{\columnwidth}%
  \def\@tempa{#1}%
  \ifx\@tempa\empty\else%
    \textbf{#1}\hspace{0.5em}%
  \fi%
}{%
  \end{boxedminipage}%
}

% header for each page of test
% be sure to change everything to what it should be (including total number of pages!)
\pagestyle{fancyplain}
\lhead{MATH 2630-080}
\chead{Test 1}
\rhead{Page \thepage \ of 9}
\cfoot{}


\begin{document}
\thispagestyle{empty}
% front page heading
\begin{Center}
    \Large{\textbf{MATH 2630-090: TEST 1 -- Summer 2024}} \\
    \large{
    June 4, 2024 \\
    Owen Henderschedt
    }
\end{Center}

\vspace*{1cm}

% space for name
\begin{FlushRight}
    Name (\textit{please print}): \underline{\hspace{12.5cm}}
\end{FlushRight}

\vspace*{1cm}

% instructions
\textbf{Instructions:}
Please read the following instructions carefully.
\begin{itemize}
    \item \textbf{Read all of the directions} for each problem.

    \item You have \textbf{75 minutes} to finish the exam.

    \item \textbf{Show work to receive credit.}
    Answers that are wrong, but still have work might receive partial credit.
    If you provide the answer to a computation without work, I will not give partial credit.

    \item Provide \textbf{exact} solutions.

    \item Do \textbf{not} write in the table below.
\end{itemize}

\vspace*{1cm}

% scoring table
\begin{center}
\begin{tabular}{|*{3}{>{\centering\arraybackslash}p{.1\linewidth}|}}
    \hline
    \textbf{Problem} & \textbf{Points} & \textbf{Score} \\
    \hline
    1 & 10 &  \\
    \hline
    2 & 12 &  \\
    \hline
    3 & 10 &  \\
    \hline
    4 & 15 &  \\
    \hline
    5 & 10 &  \\
    \hline
    6 & 12 &  \\
    \hline
    7 & 15 &  \\
    \hline
    8 & 15 &  \\
    \hline \hline
    \textbf{Total} & 99 +1 &  \\
    \hline
\end{tabular}
\end{center}


% start of test
%%%%% PROBLEM 1 %%%%%
\newpage
\begin{problem}[1. (8 points)]
Find the domain of the vector function. 

\begin{equation*}
    \vec{r}(t) = \left\langle \frac{t}{\sqrt{t-3}}, \frac{1}{(t+6)(t-5)}, \ln(t+1)\right\rangle
\end{equation*}
\end{problem}


%%%%% PROBLEM 2 %%%%%
\newpage
\begin{problem}[2. (12 points, 3 each)]
Use the vectors $\vec{v} = 4\hat{i}-3\hat{j}+\hat{k}$ and $\vec{w} = -2\hat{i}-\hat{j}+2\hat{k}$ to answer the following.

\vspace{.5cm}

(a) $\Vec{v}\cdot \vec{w}$

\vspace{.125cm}

(b) $|\vec{v}-\vec{w}|$

\vspace{.125cm}

(c) A \emph{unit} vector orthogonal to both $\vec{v}$ and $\vec{w}$.

\vspace{.125cm}

(d) The angle $\theta$ between $\vec{v}$ and $\vec{w}$ (Round to 2 decimals).

\end{problem}



%%%%% PROBLEM 3 %%%%%
\newpage
\begin{problem}[3. (10 points)]
Find the parametric equation for a line which passes through points $P(1,-1,3)$ and $Q(-2,0,7)$.
\end{problem}


%%%%% PROBLEM 4 %%%%%
\newpage
\begin{problem}[4. (15 points, 3,7,5 resp.)]
Use the two lines $L_1$ and $L_2$ to answer the following.

\begin{equation*}
    L_1: \frac{x_1-3}{-2}=\frac{y_1+1}{-4}=\frac{z_1}{2} \hspace{1cm} L_2: x_2 = 3-6t_2, y_2 = 1 + 2t_2, z_2 = -1 - 4t_2. 
\end{equation*}

\vspace{.125cm}

(a) Are $L_1$ and $L_2$ parallel? Justify your answer.

\vspace{.125cm}

(b) Are $L_1$ and $L_2$ skew lines or do they intersect? If they intersect, at which point?

\vspace{.125cm}

(c) Where does $L_1$ cross the $xz$-plane?
\end{problem}


%%%%% PROBLEM 5 %%%%%
\newpage
\begin{problem}[5. (10 points)]
Find the equation for the plane which contains points $P(-5,2,0)$ and $Q(1,1,1)$ and is perpendicular to the plane $x-y+2z = 5$.
\end{problem}


%%%%% PROBLEM 6 %%%%%
\newpage
\begin{problem}[6. (12 points)]
Evaluate the limit of the vector function, $\lim_{t\to 1}\vec{r}(t)$.

\begin{equation*}
\vec{r}(t) = (te^{-t}) \hat{i} + \left(\frac{5t^2-5}{t-1}\right) \hat{j} + \left(\cos(t-1)-\sin(t-1)\right)\hat{k}.
\end{equation*}
\end{problem}


%%%%% PROBLEM 7 %%%%%
\newpage
\begin{problem}[7. (15 points)]
Find the equation of the tangent line to the space curve given by the vector function $\vec{r}(t)$ at the time $t = 4$.

\begin{equation*}
    \vec{r}(t) = \left\langle \sqrt{t}, t+2, t^2 \right\rangle.
\end{equation*}
\end{problem}


%%%%% PROBLEM 8 %%%%%
\newpage
\begin{problem}[8. (15 points)]
If $\vec{r'}(t) = \sin(t) \hat{i} + t \hat{j} + \hat{k}$ and $\vec{r}(0) = -\hat{i} + 3\hat{k}$. Find $\vec{r}(t)$.
\end{problem}




\end{document}